\section{Discussion}
\label{sec:Diskussion}

In this section the experimental results are discussed. 

\subsection{Relative deviation from comparitive values}
First the relative deviation from comparitive values is calculated. 
The latter can be found in \cite{muccillo_1974}.
For the activation energy of strontium doped potassium bromide $W = \qty{0.66(1)}{\electronvolt}$ is used.
For the characteristic relaxation time the comparitive value is $\tau_0 = \qty{4(2)e-14}{\second}$.
Then the relative deviation is calculated with \autoref{eqn:relabw}.
The deviations for all results obtained from the two data sets using both methodes can be found in \autoref{tab:reldev}.
It stands out, that the deviations for the characteristic relaxation times are very high while those for the activation energy are quite yet aceptable.
\begin{table}
    \centering
    \caption{Relative deviation of the experimental results from comparitive values in \cite{muccillo_1974} for the activation energy and the characteristic relaxation time.}
    \label{tab:reldev}
    \sisetup{table-format=1.1, per-mode=reciprocal}
    \begin{tblr}{c c c c c}
        \toprule
        & \SetCell[c=2]{c} method 1 & & \SetCell[c=2]{c} method 2 & \\
        $b$ / $[\unit{\kelvin\per\minute}]$ & $\increment W_{\mathrm{rel}}$ / $[\unit{\percent}]$ & $\increment \tau_{0,\mathrm{rel}}$ / $[\unit{\percent}]$ & $\increment W_{\mathrm{rel}}$ / $[\unit{\percent}]$ & $\increment \tau_{0,\mathrm{rel}}$ / $[\unit{\percent}]$ \\
        \midrule
        $2$   & 33.2 & 92.6 & 34.4 & 99.9 \\
        $1.5$ & 40.4 & 96.2 & 34.0 & 99.9 \\
        \bottomrule
    \end{tblr}
\end{table}
\FloatBarrier
\subsection{Discussion for the first method}
First it has to be said, that fitting the background of both data sets was quite hard. 
The data points after the peak can not be fitted exponentially with the data points before the peak easily.
Therfore the background fit could lead to errors in the further analisys. 
\\
Likewise fitting the natural logarythm of the depolarisation current to the inverse of the temperature was just as hard, because the linear decline is not clearly visible.
So the data point are deveating from the theoretical expectaion. 
\\
Both of these could explain the deviations of the experimental results for the first method. 
The deviation of the characteristic relaxation time could be higher because the error for the activation energy propagates and a possible error of the temperature at which the current at its maximum adds. 


\subsection{Discussion for the second method}
The possible error of the bachground fit could also propagate in the analisys with the second method. 
\\
Here the linear dependence of the natural logarythm of the relaxation time is easier to regognice, but still not perfectly clear. 
But the fit values vary depending on which data points are used for the fit and which are not. 
Especially the y-axis intercept strongly varies and has a very large unvertainty. 
Therefor this explains the deviation of the experimental results.
It also explains, why the results for the characteristic relaxation times deviate so much, since they are determined from the y-axis intercept. 


\subsection{Discussion of general sources of error}
The deviation of the results for both methodes can be explained by the following general sources of error.
\\
During the conduction of the experiment it was already apparent that the experimental data is way smaler than expected. 
This could be a hint that some of the data is deveating and therefor explain, why the linear fits where not clearly visible. 
In addition an unexpected decline of the current at the end of the measurement series for larger temperatures was noticed, although a decline of the background is not expected at this point. 
This could explain in particular, why the background was so hard to fit. 
\\
